\documentclass[12pt]{article}
\usepackage[utf8]{inputenc}
\usepackage{amsmath,amssymb,amsthm}
\usepackage[margin=1in]{geometry}

\newtheorem{theorem}{Theorem}
\newtheorem{lemma}[theorem]{Lemma}
\newtheorem{corollary}[theorem]{Corollary}

\title{Problem 332: Spherical Triangles}
\author{Project Euler Solution}
\date{}

\begin{document}
\maketitle

\section{Problem Statement}

Find the sum of the areas of all spherical triangles on a unit sphere whose three sides have integer arc-lengths and whose perimeter does not exceed $100$.

\section{Mathematical Foundation}

\begin{theorem}[Girard's Theorem]
The area of a spherical triangle on a unit sphere with interior angles $\alpha, \beta, \gamma$ is
\[
A = \alpha + \beta + \gamma - \pi.
\]
\end{theorem}

\begin{proof}
Each pair of great circles determines a lune. The lune with dihedral angle $\alpha$ has area $2\alpha$. The three lunes cover the sphere ($4\pi$), covering the triangle and its antipodal image three times each, and the rest once. Thus $2(2\alpha + 2\beta + 2\gamma) = 4\pi + 4A$, yielding $A = \alpha + \beta + \gamma - \pi$.
\end{proof}

\begin{theorem}[Spherical Law of Cosines]
For sides $a, b, c$ opposite angles $\alpha, \beta, \gamma$:
\[
\cos a = \cos b \cos c + \sin b \sin c \cos \alpha.
\]
\end{theorem}

\begin{proof}
Place vertex $C$ at the north pole. Express $A$ and $B$ in Cartesian coordinates on the unit sphere at angular distances $b$ and $a$ from $C$, separated by angle $\gamma$. The dot product $\vec{OA} \cdot \vec{OB}$ gives $\cos c = \cos a \cos b + \sin a \sin b \cos \gamma$. Relabeling yields the result.
\end{proof}

\begin{lemma}[L'Huilier's Formula]
With $s = (a+b+c)/2$, the spherical excess $E = A$ satisfies
\[
\tan\frac{E}{4} = \sqrt{\tan\frac{s}{2}\,\tan\frac{s-a}{2}\,\tan\frac{s-b}{2}\,\tan\frac{s-c}{2}}.
\]
\end{lemma}

\begin{proof}
From the spherical law of cosines, express each half-angle $\alpha/2$ in terms of the sides. Using the identity $\tan^2(\alpha/2) = \frac{1-\cos\alpha}{1+\cos\alpha}$ and the product formula for $\tan(E/4)$ derived from the half-angle expressions, algebraic manipulation yields the stated product of tangent factors.
\end{proof}

\begin{lemma}[Validity Conditions]
A spherical triangle with sides $a, b, c$ exists on a unit sphere iff:
\begin{enumerate}
    \item $a + b > c$, $b + c > a$, $c + a > b$,
    \item $a + b + c < 2\pi$,
    \item $0 < a, b, c < \pi$.
\end{enumerate}
\end{lemma}

\begin{proof}
Conditions (1) and (3) ensure the arcs form a closed figure. Condition (2) prevents degeneration: if the perimeter equals or exceeds $2\pi$, the three arcs cannot enclose a proper region.
\end{proof}

\begin{corollary}
Since $\pi \approx 3.14159$, valid integer side lengths are $a, b, c \in \{1, 2, 3\}$. The perimeter bound of $100$ is automatically satisfied.
\end{corollary}

\section{Algorithm}

Enumerate all unordered triples $(a,b,c)$ with $1 \le a \le b \le c \le 3$ satisfying the triangle inequality and $a+b+c < 2\pi$. For each, compute the area via L'Huilier's formula, multiply by the number of ordered permutations, and sum.

\section{Complexity Analysis}

\begin{itemize}
    \item \textbf{Time:} $O(1)$ --- only finitely many (at most $10$) triples to check.
    \item \textbf{Space:} $O(1)$.
\end{itemize}

\section{Answer}

\[
\boxed{2717.751525}
\]

\end{document}
